\jo{Jesus will help in this section}

\textbf{Carbon scaler}\cite{CarbonScaler}: Cloud platforms are increasing their emphasis on sustainability and reducing their operational carbon footprint. A common approach for reducing carbon emissions is to exploit the temporal flexibility inherent to many cloud workloads by executing them in periods with the greenest energy and suspending them at other times. Since such suspend-resume approaches can incur long delays in job completion times, we present a new approach that exploits the elasticity of batch workloads in the cloud to optimize their carbon emissions. Our approach is based on the notion of "carbon scaling," similar to cloud autoscaling, where a job dynamically varies its server allocation based on fluctuations in the carbon cost of the grid's energy. We develop a greedy algorithm for minimizing a job's carbon emissions via carbon scaling that is based on the well-known problem of marginal resource allocation.
\\

\textbf{Machine Learning and Analytical Power Consumption Models for 5G Base Stations\cite{Piovesan2022}}: The energy consumption of the fifth generation (5G) of mobile networks is one of the major concerns of the telecom industry. However, there is not currently an accurate and tractable approach to evaluate 5G base stations' (BSs') power consumption. In this article, we propose a novel model for a realistic characterization of the power consumption of 5G multi-carrier BSs, which builds on a large data collection campaign. At first, we define a machine learning architecture that allows modeling multiple 5G BS products. Then we exploit the knowledge gathered by this framework to derive a realistic and analytically tractable power consumption model, which can help driving both theoretical analyses as well as feature standardization, development, and optimization frameworks. Notably, we demonstrate that this model has high precision, and it is able to capture the benefits of energy-saving mechanisms. We believe this analytical model represents a fundamental tool for understanding 5G BSs' power consumption and accurately optimizing the network energy efficiency.\\

\textbf{Neural Networks for Cellular Base Station Switching\cite{Donevski2019}.}: The adoption of base station sleep modes is considered one of the most effective approaches for the reduction of the energy consumption of radio access networks. Sleep modes allow the switch-off of base stations in periods of low traffic, and their successive switch-on when traffic increases. The selection of the appropriate instants to switch base stations on and off requires an accurate prediction of the traffic loads in the near future. In this paper we explore the performance of machine learning techniques for traffic prediction and for the selection of the instants when to switch off and on base stations, considering a heterogeneous network portion comprising one macro cell and six small cells within the macro cell coverage. We experiment with two machine learning approaches. The first aims at short-term traffic load estimation, and from this derives the best combination of switching decisions. The second performs both traffic estimation and switching optimization at one time. For both approaches we develop artificial neural network implementations based on the Dense Neural Network and Recurrent Neural Network paradigms. Testing the two approaches on real traffic data, we observe very good performance in terms of both quality of service and energy saving.\\


\textbf{Modeling User Transfer During Dynamic Carrier Shutdown in Green 5G Networks\cite{user_transfer_5g}.}: The energy consumption of the fifth generation (5G) of cellular technology is concerning for the mobile industry and the entire society. To minimize the environmental footprint and economic costs of 5G, it is necessary to adapt the transmission capabilities of networks to end-users’ quality of service requirements. In this paper, we focus on the carrier shutdown approach that enables a base station (BS) to autonomously switch off during low traffic periods, by transferring its load to neighbouring active BSs. More specifically, we propose a data-driven framework, constructed through real network measurements, which statistically characterizes the user equipment (UE) transfer across neighbouring BSs, when carrier shutdown operates. The implementation of this framework allows the 5G system to determine a poor load distribution due to energy saving mechanisms, prevent drastic reductions in UE performance, and ultimately estimate energy savings when activating carrier shutdown.\\


\textbf{Auric: Using Data-driven Recommendation to Automatically Generate Cellular Configuration \cite{Mahimkar2021}}: Cellular service providers add carriers in the network in order to support the increasing demand in voice and data traffic and provide good quality of service to the users. Addition of new carriers requires the network operators to accurately configure their parameters for the desired behaviors. This is a challenging problem because of the large number of parameters related to various functions like user mobility, interference management and load balancing. Furthermore, the same parameters can have varying values across different locations to manage user and traffic behaviors as planned and respond appropriately to different signal propagation patterns and interference. Manual configuration is time-consuming, tedious and error-prone, which could result in poor quality of service. In this paper, we propose a new data-driven recommendation approach Auric to automatically and accurately generate configuration parameters for new carriers added in cellular networks. Our approach incorporates new algorithms based on collaborative filtering and geographical proximity to automatically determine similarity across existing carriers. We conduct a thorough evaluation using real-world LTE network data and observe a high accuracy (96\%) across a large number of carriers and configuration parameters. We also share experiences from our deployment and use of Auric in production environments.\\

\textbf{A Survey on 5G Radio Access Network Energy Efficiency: Massive MIMO, Lean Carrier Design, Sleep Modes, and Machine Learning\cite{survey_5g_energy_efficiency}}:Cellular networks have changed the world we are living in, and the fifth generation (5G) of radio technology is expected to further revolutionise our everyday lives by enabling a high degree of automation, through its larger capacity, massive connectivity, and ultra-reliable low-latency communications. In addition, the third generation partnership project (3GPP) new radio (NR) specification also provides tools to significantly decrease the energy consumption and the green house emissions of next generations networks, thus contributing towards information and communication technology (ICT) sustainability targets. In this survey paper, we thoroughly review the state-of-the-art on current energy efficiency research. We first categorize and carefully analyse the different power consumption models and energy efficiency metrics, which have helped to make progress on the understanding of green wireless networks. Then, as a main contribution, we survey in detail —from a theoretical and a practical viewpoint— the main energy efficiency enabling technologies that 3GPP NR provides, together with their main benefits and challenges. Special attention is paid to four key enabling technologies, i.e., massive multiple-input multiple-output (MIMO), lean carrier design, and advanced idle modes, together with the role of artificial intelligence capabilities. We dive into their implementation and operational details and thoroughly discuss their optimal operation points and theoretical trade offs from an energy consumption perspective. This will help the reader to grasp the fundamentals of —and the status on— green wireless networking. Finally, the areas of research where more effort is needed to make future wireless networks greener are also discussed.\\
